\documentclass[11pt]{article}
\usepackage[margin=1in]{geometry}
\usepackage{amsmath,amssymb,amsthm,hyperref}
\newtheorem{theorem}{Theorem}
\newtheorem{lemma}{Lemma}
\newcommand{\PN}{\mathcal P_N}
\title{Nonexistence of the proposed persistence density}
\author{Principia Math In Progress}
\date{28 September 2026}
\begin{document}
\maketitle
\begin{abstract}
For every fixed integer sample depth $N\ge3$, the persistent integer
bases defined by Hyman's continued-fraction criterion have no natural
density. Thus the proposed density $1/12$ does not exist. The proof
uses two open intervals of fractional logarithms, one entirely
persistent and one entirely nonpersistent, and their repetition in
every decimal decade.
\end{abstract}

\section*{Task 1: the precise density question}
Fix $N\ge3$ and write $\log$ for the natural logarithm. For an integer
$b\ge2$ with irrational $\alpha_b=\log_{10}b$, write
$\alpha_b=[a_0;a_1,a_2,\ldots]$, and let $q_k$ be its convergent
denominators, with $q_{-1}=0$, $q_0=1$. Following Corollary 3.16
of \cite{H}, define
\[
 \PN=\{b\ge2:\text{ some }k\ge0\text{ has }q_k\le N
                    \text{ and }a_{k+1}q_k>N\log N\}.
\]
The rational-logarithm bases may be included or excluded arbitrarily;
they have density zero, as shown below. Conjecture 3.17 proposes
\[
 \lim_{B\to\infty}\frac{\#(\PN\cap[2,B])}{B-1}=\frac1{12}.
\]
The sample depth is retained explicitly here: the source defines
persistence at depth $N$ and surveys $N=10{,}000$, with an extended
survey at $N=200{,}000$. Its supplementary Section 9.8 also keeps $N$
fixed in the continued-fraction cutoff while increasing the base bound.

\begin{theorem}
For every fixed integer $N\ge3$, the natural density of $\PN$ does not
exist. In particular, Conjecture 3.17 fails for every such choice of
the fixed sample depth, including both survey depths.
\end{theorem}

\section*{Two intervals with opposite classifications}
Adding an integer to a continued fraction changes only $a_0$.
The recurrence $q_k=a_kq_{k-1}+q_{k-2}$ therefore shows that the
persistence criterion depends only on
$\theta=\{\log_{10}b\}\in(0,1)$.

First choose $A=N^2$. Since $\log N<N$, we have $A>N\log N$.
For every irrational
\[
 \theta\in J_+:=\left(\frac{A}{A+1},1\right),
\]
one has $a_1=1$, $q_1=1$, and
\[
 a_2=\left\lfloor\frac{\theta}{1-\theta}\right\rfloor
       \ge A>N\log N.
\]
Thus every such $\theta$ meets the persistence criterion already at
$k=1$. This argument does not require admitting the index $k=0$.

Next consider a continued-fraction prefix consisting entirely of ones.
Let
\[
 p_{-1}=1,\quad p_0=0,\quad q_{-1}=0,\quad q_0=1,
 \qquad p_j=p_{j-1}+p_{j-2},\quad q_j=q_{j-1}+q_{j-2}.
\]
Choose $L$ with $q_L>N$. The values
\[
 [0;\underbrace{1,\ldots,1}_{L\text{ entries}},t]
       =\frac{p_Lt+p_{L-1}}{q_Lt+q_{L-1}},\qquad 1<t<\infty,
\]
fill the open interval $J_-$ between
\[
 u=\frac{p_L}{q_L},\qquad
 v=\frac{p_L+p_{L-1}}{q_L+q_{L-1}}.
\]
The determinant $p_Lq_{L-1}-p_{L-1}q_L=(-1)^{L-1}$,
proved immediately by induction, gives
\[
 |u-v|=\frac1{q_L(q_L+q_{L-1})}>0.
\]
Every irrational number in $J_-$ has $a_1=\cdots=a_L=1$.
Since its $L$th denominator exceeds $N$, every index with $q_k\le N$
has $k<L$. Consequently
\[
 a_{k+1}q_k=q_k\le N<N\log N,
\]
where $\log N>1$ for $N\ge3$. None of these numbers is persistent.

\section*{Decimal intervals contradict a natural density}
If an integer $b\ge2$ has $\log_{10}b=p/q$ in lowest terms, then
$b^q=10^p$. Taking the exponent of the prime $2$ gives
$q\,v_2(b)=p$, so $q=1$. Such bases are exactly the integral powers
of ten. There are only $O(\log B)$ of them up to $B$.

Put
\[
 c_+=10^{A/(A+1)},\quad d_+=10,\qquad
 c_-=10^{\min(u,v)},\quad d_-=10^{\max(u,v)}.
\]
For every nonnegative integer $m$, all integer bases in
$(c_+10^m,d_+10^m)$ are persistent, while none in
$(c_-10^m,d_-10^m)$ is persistent. These open intervals lie within
a single decimal decade, so contain no powers of ten.

Suppose a natural density $\delta$ existed and set
$C(B)=\#(\PN\cap[2,B])$. Equivalently, $C(B)=\delta B+o(B)$.
For fixed $1\le c<d\le10$, subtracting this asymptotic at the two
endpoints gives
\begin{equation}\label{counts}
 \#\bigl(\PN\cap(c10^m,d10^m)\bigr)
    =\delta(d-c)10^m+o(10^m).
\end{equation}
The possible integer endpoints change this equation by at most two.
For $(c,d)=(c_+,d_+)$ the left side is
$(d_+-c_+)10^m+O(1)$, so \eqref{counts} forces $\delta=1$.
For $(c,d)=(c_-,d_-)$ it is zero, forcing $\delta=0$.
This contradiction proves the theorem.

This is a refutation of the ordinary base-counting limit stated in
the source. A logarithmically weighted density or a joint limit with
$N=N(B)$ would be a different mathematical question. The proof uses
the source's exact continued-fraction definition and does not require
any claimed CMI estimates or survey classifications.

\section*{Reproducible checks}
The accompanying exact verifier checks the continued-fraction
recurrences, cylinder endpoint identities, and all relevant
denominator cutoffs for every integer $3\le N\le10{,}000$, and
also for $N=200{,}000$ and $N=10^9$. It verifies
$1<\log3$ by a rational lower Taylor bound and uses
$\log N\le N-1<N$ analytically. At $N=3$, the inequalities
$9^{10}>10^9$, $9<10$, $4^5>10^3$, and $4^8<10^5$
place bases $9$ and $4$ in the two respective cylinders without
floating-point logarithms. These checks supplement the universal proof.

\section*{Proposed follow-up tasks}
\begin{description}
\item[Task 2.] Determine the lower and upper natural densities of
$\PN$ at the survey depth $N=10{,}000$.
\item[Task 3.] Determine whether $\PN$ has a logarithmic density and
express it in terms of continued-fraction cylinders.
\item[Task 4.] Determine the asymptotic behavior, as $N\to\infty$,
of the Lebesgue measure of persistent rotation numbers in $(0,1)$.
\item[Task 5.] Identify growth regimes $N=N(B)$ for which the
persistence proportion up to base $B$ has a limit, and compute it.
\end{description}

\begin{thebibliography}{9}
\bibitem{H} J. M. Hyman (Mac Hyman),
\emph{Why Eight Percent of Benford Sequences Never Converge},
\href{https://arxiv.org/abs/2603.18243}{arXiv:2603.18243v1} (2026).
Definition 3.11, Corollary 3.16 and Conjecture 3.17, pp.~16--18;
Supplementary Information, Section 9.8, pp.~33--34,
\href{https://github.com/machyman/hyman2026eight/tree/0ffdd7cf0e6f660c789b5fe95bdcef61fa1d48ca/paper}{author repository}.
\end{thebibliography}
\end{document}
